\begin{ouverture}[college1]{6}{Reconnaître une transformation chimique}{Partie 1 · Organisation et transformations de la matière}{Comment prouver qu’une nouvelle espèce s’est formée ?}
\objectif{Distinguer mélange, changement d’état et transformation chimique.}
\objectif{Exploiter une banque de tests et un témoin.}
\objectif{Identifier réactifs et produits et lire une équation fournie.}
\prerequis{Formules chimiques et états physiques.}
\end{ouverture}
\pcPage{Le cours}{Reconnaître une transformation chimique}
\phantomsection\label{ess-6}
\essSection{1}{Observer puis interpréter}
Une transformation chimique produit de nouvelles espèces à partir d’espèces consommées. Les réactifs sont consommés ; les produits se forment. Un changement de couleur, une température qui varie ou des bulles sont des indices, pas des preuves isolées. Une ébullition ou un dégazage produit aussi des bulles sans créer nécessairement une nouvelle espèce. Il faut confronter plusieurs observations et des tests.\par
\pcFigure{figures/c06-tests.pdf}{Banque de tests simplifiée : les essais avec flamme sont exclusivement réalisés par le professeur sur de très petites quantités.}
\essSection{2}{Utiliser des tests}
Dans la banque fournie : le dioxyde de carbone trouble l’eau de chaux ; le dioxygène ravive une bûchette incandescente ; un très petit prélèvement de dihydrogène donne une petite détonation avec une flamme, uniquement sous conduite du professeur. Le sulfate de cuivre anhydre blanc bleuit en présence d’eau. Un test négatif peut aussi venir d’une quantité trop faible ou d’une mauvaise manipulation.\par
\clearpage
\essSection{3}{Des atomes qui se réorganisent}
Lors d’une réaction chimique, les atomes se réarrangent : leur nature et leur nombre se conservent dans le modèle étudié. Les molécules peuvent disparaître et d’autres apparaître. Une équation fournie résume le bilan ; la flèche indique le sens de la transformation, pas une égalité entre masses de chaque espèce. Les coefficients équilibrent les inventaires sans modifier les formules.\par
\pcFigure{figures/c06-reaction.pdf}{$2\,\ce{H2}+\ce{O2}\rightarrow2\,\ce{H2O}$ : mêmes atomes avant et après, assemblages différents.}
\essSection{4}{Varier les situations}
Une réaction peut être rapide ou lente. La corrosion d’un métal est une transformation chimique : le métal réagit avec son environnement et forme d’autres espèces. Une réaction entre un acide et un métal peut produire du dihydrogène ; ce n’est pas vrai pour tous les couples métal-solution. Les réactions acide-base seront approfondies en 3e. Une banque documentaire permet déjà d’identifier réactifs, produits et précautions sans réaliser d’essai improvisé.\par
\begin{pcEncadre}{method}{Un exemple expliqué}On fait réagir, en classe, un matériau avec une solution autorisée. Le gaz recueilli trouble l’eau de chaux, contrairement à un prélèvement témoin. La banque permet de conclure à la présence de dioxyde de carbone dans les conditions du test. Le trouble ne donne pas directement sa quantité ni la liste de toutes les espèces présentes.\end{pcEncadre}
\pcPage{Expérimenter et retenir}{Reconnaître une transformation chimique}
\phantomsection\label{activite-6}
\begin{pcEncadre}{method}{Investigation en classe}\textbf{Question.} Une effervescence correspond-elle à une création de gaz ? Le professeur compare un dégazage d’eau gazeuse et une réaction autorisée entre un carbonate et un acide dilué. Utiliser un témoin et la même banque de tests. Dresser deux colonnes : « faits observés » et « interprétation ». Ne jamais fermer une réaction produisant un gaz dans un récipient rigide improvisé. \textbf{À rendre :} identification justifiée du gaz et différence entre séparation d’un gaz déjà présent et formation chimique.\end{pcEncadre}
\pcSousTitre{Ce que je dois savoir expliquer}\begin{questions}
\item Un réactif est consommé et un produit se forme.
\item Un indice visuel ne prouve pas à lui seul une réaction.
\item Un témoin aide à interpréter un test.
\item Les atomes se conservent ; les molécules peuvent changer.
\item On lit une équation en comptant les atomes de chaque élément.
\end{questions}
\begin{pcEncadre}{attention}{Une erreur à éviter}Changer $\ce{H2O}$ en $\ce{H2O2}$ pour équilibrer une équation change le produit : il faut agir sur les coefficients, pas sur les indices.\end{pcEncadre}
\pcRepere{cor-activite-6}{Pistes pour l’investigation}\pcGrille
\clearpage\markright{Je démarre}
\rubrique{Je démarre}
\par\noindent\begin{minipage}{\linewidth}
\begin{exercice}[1]{Un changement physique}
\begin{questions}
\item La glace fond. Une nouvelle espèce chimique apparaît-elle ?
\item Comment classer cette transformation ?
\end{questions}
\end{exercice}
\end{minipage}
\par\noindent\begin{minipage}{\linewidth}
\begin{exercice}[1]{Réactif ou produit}
\begin{questions}
\item Dans une réaction, comment nommer une espèce consommée ?
\item Et une espèce qui se forme ?
\end{questions}
\end{exercice}
\end{minipage}
\par\noindent\begin{minipage}{\linewidth}
\begin{exercice}[1]{Le trouble observé}
\begin{questions}
\item Un gaz trouble l’eau de chaux. Quelle espèce la banque identifie-t-elle ?
\item Ce résultat mesure-t-il directement une masse ?
\end{questions}
\end{exercice}
\end{minipage}
\par\noindent\begin{minipage}{\linewidth}
\begin{exercice}[1]{Le fer rouille}
\begin{questions}
\item La corrosion du fer est-elle un simple changement d’état ?
\item Que peut-on affirmer sur les espèces ?
\end{questions}
\end{exercice}
\end{minipage}
\par\noindent\begin{minipage}{\linewidth}
\begin{exercice}[1]{Les bulles ne suffisent pas}
\begin{questions}
\item Une eau bout. Les bulles prouvent-elles une réaction chimique ?
\item Quelle autre explication convient ?
\end{questions}
\end{exercice}
\end{minipage}
\clearpage\markright{Je m’entraîne}
\rubrique{Je m’entraîne}
\par\noindent\begin{minipage}{\linewidth}
\begin{exercice}[2]{Lire un bilan}
\begin{questions}
\item Dans $\ce{C + O2 -> CO2}$, nommer réactifs et produit.
\item Combien d’atomes O sont comptés de chaque côté ?
\end{questions}
\end{exercice}
\end{minipage}
\par\noindent\begin{minipage}{\linewidth}
\begin{exercice}[2]{Un coefficient utile}
\begin{questions}
\item Lire $2\,\ce{H2}+\ce{O2}\rightarrow2\,\ce{H2O}$. Compter H et O avant et après.
\item Combien de molécules d’eau sont représentées ?
\end{questions}
\end{exercice}
\end{minipage}
\par\noindent\begin{minipage}{\linewidth}
\begin{exercice}[2]{Le témoin}
\pcFigureEx{figures/ex-temoin.pdf}{Un seul tube reçoit le gaz étudié ; les autres conditions sont comparables.}
\begin{questions}
\item Deux tubes reçoivent la même eau de chaux ; seul l’un reçoit le gaz étudié. Pourquoi conserver l’autre ?
\item Que conclure si les deux troublent avant tout ajout ?
\end{questions}
\end{exercice}
\end{minipage}
\par\noindent\begin{minipage}{\linewidth}
\begin{exercice}[2]{Une identification limitée}
\begin{questions}
\item Un test de dioxyde de carbone est positif dans un mélange de gaz. Peut-on affirmer que le mélange est pur ?
\item Comment formuler la conclusion ?
\end{questions}
\end{exercice}
\end{minipage}
\par\noindent\begin{minipage}{\linewidth}
\begin{exercice}[2]{Un test négatif}
\begin{questions}
\item Le sulfate de cuivre anhydre ne bleuit pas. Peut-on conclure qu’aucune trace d’eau n’existe ?
\item Donner deux vérifications.
\end{questions}
\end{exercice}
\end{minipage}
\clearpage\markright{Je raisonne}
\rubrique{Je raisonne}
\par\noindent\begin{minipage}{\linewidth}
\begin{exercice}[3]{Corriger un équilibrage}
\begin{questions}
\item Un élève transforme $\ce{H2 + O2 -> H2O}$ en $\ce{H2 + O2 -> H2O2}$. Pourquoi sa démarche est-elle incorrecte ?
\item Donner le bilan correct pour former de l’eau.
\end{questions}
\end{exercice}
\end{minipage}
\par\noindent\begin{minipage}{\linewidth}
\begin{exercice}[3]{Séparer les faits des idées}
\begin{questions}
\item Classer : « des bulles apparaissent » et « du dioxyde de carbone se forme ».
\item Quelle démarche relie la première phrase à la seconde ?
\end{questions}
\end{exercice}
\end{minipage}
\par\noindent\begin{minipage}{\linewidth}
\begin{exercice}[3]{Un métal et un acide}
\begin{questions}
\item Une banque indique : « zinc + solution acide : formation de dihydrogène ». Pourquoi ne pas généraliser à tous les métaux ?
\item Peut-on goûter le liquide pour vérifier son acidité ?
\end{questions}
\end{exercice}
\end{minipage}
\par\noindent\begin{minipage}{\linewidth}
\begin{exercice}[3]{Un gaz déjà présent}
\begin{questions}
\item Deux expériences libèrent le même gaz, l’une par dégazage et l’autre par réaction. Le test final distingue-t-il à lui seul les mécanismes ?
\item Quelle information faut-il ajouter ?
\end{questions}
\end{exercice}
\end{minipage}
\par\noindent\begin{minipage}{\linewidth}
\begin{exercice}[3]{Concevoir une preuve}
\begin{questions}
\item On veut montrer qu’une poudre réagit et ne se contente pas de se dissoudre. Proposer deux recherches.
\item Pourquoi une simple disparition visuelle est-elle insuffisante ?
\end{questions}
\end{exercice}
\end{minipage}
\pcPage{Défi du chapitre}{Reconnaître une transformation chimique}
\phantomsection\label{defi-ex-6}
\pcDefiTitre{6}{Deux flacons effervescents}
\begin{questions}
\item Le flacon A contient une boisson gazeuse ouverte ; B met en contact un carbonate et un acide autorisés. Les deux gaz troublent l’eau de chaux. Que montrent les tests ?
\item Pourquoi ne prouvent-ils pas deux réactions identiques ?
\item Proposer le rôle d’un témoin et une conclusion rédigée.
\end{questions}
\pcGrille
